\chapter{Nonatomic measures, Lyapunov's theorem and continuum economies}
\label{chap:lyapunov}

Economies with a continuum of agents (Aumann 1964) model perfect competition: each agent is negligible. The key
mathematical tool is Lyapunov's (1940) theorem that the range of a finite-dimensional nonatomic vector measure is
convex. This chapter builds it from Sierpiński's intermediate value theorem for nonatomic measures.

\section{Sierpiński's theorem}
\label{sec:sierpinski}

\begin{definition}[Nonatomic measure]
\label{def:nonatomic}
\lean{Econ.Lyapunov.Nonatomic, Econ.Lyapunov.measureReal_pos_iff, Econ.Lyapunov.Nonatomic.real}
\leanok
A measure $\mu$ is nonatomic if every measurable set of x

\begin{lemma}[Small subsets]
\label{lem:small_subsets}
\lean{Econ.Lyapunov.Nonatomic.exists_subset_half, Econ.Lyapunov.exists_subset_measureReal_pos_le}
\leanok
\uses{def:nonatomic}
For a finite nonatomic $\mu$, every set of positive measure has a subset of positive measure at most half its measure,
hence subsets of arbitrarily small positive measure.
\end{lemma}
\begin{proof}
\leanok
If $0 < \mu(u) < \mu(s)$, one of $u$ and $s \setminus u$ has positive measure at most $\mu(s)/2$. Iterate $k$ times.
\end{proof}

\begin{theorem}[Sierpiński 1922]
\label{thm:sierpinski}
\lean{Econ.Lyapunov.exists_subset_measureReal_eq}
\leanok
\uses{lem:small_subsets}
If $\mu$ is finite and nonatomic, $A$ measurable and $0 \le c \le \mu(A)$, there is a measurable $B \subseteq A$ with
$\mu(B) = c$.
\end{theorem}
\begin{proof}
\leanok
Greedy exhaustion: at each stage add a subset $D_k$ of $A \setminus B_k$ with $\mu(B_k) + \mu(D_k) \le c$ and at least half
the supremum $S_k$ of such sizes. The measures $\mu(B_k)$ increase to $\mu(\bigcup B_k) \le c$, so $\mu(D_k) \to 0$ and
$S_k \to 0$. If $\mu(\bigcup B_k) < c$, a small positive subset of $A \setminus \bigcup B_k$ would be admissible at every
stage, contradicting $S_k \to 0$.
\end{proof}

\begin{corollary}[Fine partitions]
\label{cor:fine_partition}
\lean{Econ.Lyapunov.exists_partition_measureReal_le_aux, Econ.Lyapunov.exists_partition_measureReal_le}
\leanok
\uses{thm:sierpinski}
Every measurable set of a finite nonatomic measure space splits into finitely many disjoint measurable pieces of measure
at most $\delta$.
\end{corollary}
\begin{proof}
\leanok
Split off a piece of measure exactly $\delta$ and induct.
\end{proof}

\begin{lemma}[Sums of nonatomic measures]
\label{lem:nonatomic_sum}
\lean{Econ.Lyapunov.nonatomic_sum}
\leanok
\uses{thm:sierpinski}
A finite sum of finite nonatomic measures is nonatomic.
\end{lemma}
\begin{proof}
\leanok
Some summand $\mu_i$ charges $s$; a subset $t$ with $\mu_i(t) = \mu_i(s)/2$ has positive sum-measure, and so does $s \setminus t$.
\end{proof}

\begin{lemma}[Absolute continuity preserves nonatomicity]
\label{lem:nonatomic_ac}
\lean{Econ.Lyapunov.Nonatomic.of_absolutelyContinuous}
\leanok
\uses{thm:sierpinski}
If $\lambda$ and $\nu$ are finite, $\lambda$ is nonatomic and $\nu \ll \lambda$, then $\nu$ is nonatomic. In particular
measures with densities, $B \mapsto \int_B f\,d\lambda$ for $f \ge 0$ integrable, are nonatomic.
\end{lemma}
\begin{proof}
\leanok
If $s$ were a $\nu$-atom, halve it repeatedly in $\lambda$-measure (Theorem~\ref{thm:sierpinski}), each time keeping the half
of full $\nu$-measure. The intersection is $\lambda$-null but has $\nu$-measure $\nu(s) > 0$.
\end{proof}

\section{Lyapunov's convexity theorem}
\label{sec:lyapunov}

\begin{lemma}[Few fractional weights]
\label{lem:vertex}
\lean{Econ.Lyapunov.exists_few_fractional_aux, Econ.Lyapunov.exists_few_fractional}
\leanok
Let $w_l \in \mathbb{R}^\eta$ ($\eta$ finite) for $l$ in a finite set $\iota$ and $d \in [0,1]^\iota$. There is
$d' \in [0,1]^\iota$ with $\sum_l d'_l w_l = \sum_l d_l w_l$ and at most $|\eta|$ coordinates of $d'$ in $(0,1)$.
\end{lemma}
\begin{proof}
\leanok
If more than $|\eta|$ coordinates are fractional, the corresponding $w_l$ are linearly dependent; move $d$ along a nontrivial
relation $g$ until the first fractional coordinate with $g_l \ne 0$ reaches $0$ or $1$. The fractional count drops; induct.
\end{proof}

\begin{lemma}[One rounding step]
\label{lem:lyapunov_round}
\lean{Econ.Lyapunov.FracState, Econ.Lyapunov.measureReal_le_sum, Econ.Lyapunov.FracState.round}
\leanok
\uses{lem:vertex, cor:fine_partition, lem:nonatomic_sum}
A fractional set is a measurable $B \subseteq A$ with disjoint pieces $Q_l \subseteq A \setminus B$ carrying weights
$d_l \in [0,1]$ such that $\mu(B) + \sum_l d_l \mu(Q_l) = t\mu(A)$. Any fractional set can be replaced by one with a larger
$B$, at most $|\eta|$ pieces, each of $\nu$-measure at most $\delta$, where $\nu = \sum_i \mu_i$.
\end{lemma}
\begin{proof}
\leanok
Cut each piece into pieces of $\nu$-measure at most $\delta$ (Corollary~\ref{cor:fine_partition}, $\nu$ nonatomic by
Lemma~\ref{lem:nonatomic_sum}), apply Lemma~\ref{lem:vertex} to the vectors $\mu(R_k)$, add pieces of weight $1$ to $B$ and
keep the at most $|\eta|$ fractional ones.
\end{proof}

\begin{theorem}[Lyapunov 1940]
\label{thm:lyapunov}
\lean{Econ.Lyapunov.exists_subset_measureReal_eq_smul, Econ.Lyapunov.range, Econ.Lyapunov.convex_range}
\leanok
\uses{lem:lyapunov_round}
Let $(\mu_i)_{i \in \eta}$ be finitely many finite nonatomic measures. For every measurable $A$ and $t \in [0,1]$ there is a
measurable $B \subseteq A$ with $\mu_i(B) = t\mu_i(A)$ for all $i$. Hence the range $\{(\mu_i(B))_i\}$ is convex.
\end{theorem}
\begin{proof}
\leanok
Start from $B = \emptyset$ and the single piece $A$ of weight $t$. After round $k$ (with $\delta = 1/(k+1)$) the fractional
remainder is at most $|\eta|/(k+1)$ in every coordinate, and the sets $B$ increase, so their union has measure exactly
$t\mu(A)$. For convexity, given $E, F$ and $a + b = 1$ take $(E \cap F) \cup B_1 \cup B_2$ with $B_1 \subseteq E \setminus F$
of measure $a\mu(E \setminus F)$ and $B_2 \subseteq F \setminus E$ of measure $b\mu(F \setminus E)$.
\end{proof}

\section{Lyapunov's theorem for integrals}
\label{sec:lyapunov_integral}

\begin{theorem}[Lyapunov for indefinite integrals]
\label{thm:lyapunov_integral}
\lean{Econ.Lyapunov.exists_subset_setIntegral_eq_smul, Econ.Lyapunov.setIntegral_eq_measureReal_sub,
  Econ.Lyapunov.partMeasure, Econ.Lyapunov.isFiniteMeasure_partMeasure, Econ.Lyapunov.nonatomic_partMeasure,
  Econ.Lyapunov.setIntegral_apply_eq_partMeasure}
\leanok
\uses{thm:lyapunov, lem:nonatomic_ac}
Let $\lambda$ be a finite nonatomic measure and $f\colon \Omega \to \mathbb{R}^\eta$ integrable, $\eta$ finite. For every
measurable $A$ and $t \in [0,1]$ there is a measurable $B \subseteq A$ with $\int_B f\,d\lambda = t \int_A f\,d\lambda$.
\end{theorem}
\begin{proof}
\leanok
The $2|\eta|$ measures $f_i^{\pm}\lambda$ are finite (as $f$ is integrable) and nonatomic (being absolutely continuous
with respect to $\lambda$), and $\int_B f_i = f_i^+\lambda(B) - f_i^-\lambda(B)$. Apply Theorem~\ref{thm:lyapunov}.
\end{proof}

\begin{corollary}[Convex range]
\label{cor:convex_integral_range}
\lean{Econ.Lyapunov.integralRange, Econ.Lyapunov.partDiff, Econ.Lyapunov.convex_integralRange}
\leanok
\uses{thm:lyapunov_integral}
The range $\{\int_B f\,d\lambda : B \text{ measurable}\}$ is convex.
\end{corollary}
\begin{proof}
\leanok
It is the image of the convex range of $(f_i^{\pm}\lambda)_{i}$ under the linear map $v \mapsto (v_{i,+} - v_{i,-})_i$.
\end{proof}

\begin{lemma}[Gluing]
\label{lem:gluing}
\lean{Econ.Lyapunov.exists_subset_piecewise_setIntegral_eq, Econ.Lyapunov.exists_piecewise_integral_eq}
\leanok
\uses{thm:lyapunov_integral}
For integrable $g_1, g_2$, measurable $A$ and $t \in [0,1]$, some measurable $B \subseteq A$ has
$\int_A (g_1 1_B + g_2 1_{B^c}) = t \int_A g_1 + (1-t)\int_A g_2$.
\end{lemma}
\begin{proof}
\leanok
Apply Theorem~\ref{thm:lyapunov_integral} to $(g_1, g_2)\colon \Omega \to \mathbb{R}^{\eta \sqcup \eta}$ to get
$\int_B g_k = t \int_A g_k$; then $\int_{A \setminus B} g_2 = (1-t)\int_A g_2$.
\end{proof}

\begin{theorem}[Integrals of correspondences are convex; Richter 1963, Aumann 1965]
\label{thm:convex_selection_integral}
\lean{Econ.Lyapunov.selectionIntegral, Econ.Lyapunov.convex_selectionIntegral}
\leanok
\uses{lem:gluing}
For a finite nonatomic $\lambda$ and any correspondence $F\colon \Omega \rightrightarrows \mathbb{R}^\eta$, the set
$\int F = \{\int g : g \text{ integrable}, g(a) \in F(a) \text{ a.e.}\}$ is convex.
\end{theorem}
\begin{proof}
\leanok
Given selections $g_1, g_2$ and $t$, glue them on the set $B$ of Lemma~\ref{lem:gluing}: $g_1 1_B + g_2 1_{B^c}$ is again
a selection, with integral $t\int g_1 + (1-t)\int g_2$.
\end{proof}

\begin{theorem}[Witnesses]
\label{thm:lyapunov_integral_witness}
\lean{Econ.Lyapunov.nonatomic_volume_Icc, Econ.Lyapunov.exists_half_integral_id}
\leanok
\uses{def:nonatomic, thm:lyapunov_integral}
Lebesgue measure on $[0,1]$ is nonatomic; hence some measurable $B \subseteq [0,1]$ has $\int_B a\,da = 1/4$.
\end{theorem}
\begin{proof}
\leanok
$x \mapsto \lambda(s \cap (-\infty, x])$ is $1$-Lipschitz, so by the intermediate value theorem it takes the value
$\lambda(s)/2$.
\end{proof}

\section{Exchange economies with a continuum of traders}
\label{sec:continuum_economy}

Traders form a finite measure space $(\Omega, \mu)$ (Aumann 1964); trader $a$ has preferences $\succsim_a$ on
$\mathbb{R}^\iota$ and endowment $e(a) \ge 0$, $e$ integrable. An allocation $x$ is feasible if it is integrable,
nonnegative a.e.\ and $\int x\,d\mu = \int e\,d\mu$. A coalition $S$ (measurable, $\mu(S) > 0$) blocks $x$ if some $y$ with
$\int_S y = \int_S e$ makes almost every member strictly better off; the core consists of the feasible allocations that no
coalition blocks.

\begin{theorem}[Walrasian allocations are in the core]
\label{thm:continuum_we_core}
\lean{Econ.Continuum.ContinuumEconomy, Econ.Continuum.ContinuumEconomy.Feasible, Econ.Continuum.ContinuumEconomy.Blocks, Econ.Continuum.ContinuumEconomy.InCore, Econ.Continuum.ContinuumEconomy.IsWalrasianEquilibrium, Econ.Continuum.ContinuumEconomy.costCLM, Econ.Continuum.ContinuumEconomy.costCLM_apply, Econ.Continuum.ContinuumEconomy.cost_setIntegral, Econ.Continuum.ContinuumEconomy.IsWalrasianEquilibrium.inCore}
\leanok
If $(p, x)$ is a Walrasian equilibrium (almost every $x(a)$ is in trader $a$'s demand at wealth $p \cdot e(a)$), then $x$
is in the core.
\end{theorem}
\begin{proof}
\leanok
If $S$ blocks via $y$, then $p \cdot y(a) > p \cdot e(a)$ for almost every $a \in S$, so
$p \cdot \int_S y > p \cdot \int_S e$ since $\mu(S) > 0$; but $\int_S y = \int_S e$.
\end{proof}

\begin{theorem}[A continuum witness]
\label{thm:continuum_witness}
\lean{Econ.Continuum.ContinuumEconomy.linContinuum, Econ.Continuum.ContinuumEconomy.linContinuum_isWalrasianEquilibrium, Econ.Continuum.ContinuumEconomy.linContinuum_inCore, Econ.Continuum.ContinuumEconomy.linContinuum_blocks_zero, Econ.Continuum.ContinuumEconomy.linContinuum_zero_not_inCore}
\leanok
\uses{thm:continuum_we_core}
Traders uniform on $[0,1]$, each with endowment $(1,1)$ and utility $x_0 + x_1$: autarky is a Walrasian equilibrium at
prices $(1,1)$, hence in the core, while the grand coalition blocks the zero allocation.
\end{theorem}
\begin{proof}
\leanok
Every affordable bundle has $y_0 + y_1 \le 2$.
\end{proof}


\section{Aumann's core equivalence theorem}
\label{sec:aumann_core}

Fix a core allocation $x$. Trader $a$'s \emph{improving net trades} are
$G(a) = \{z : e(a) + z \ge 0,\ e(a) + z \succ_a x(a)\} \cup \{0\}$.

\begin{lemma}[No negative aggregate improvement]
\label{lem:aumann_no_neg}
\lean{Econ.Continuum.ContinuumEconomy.improve, Econ.Continuum.ContinuumEconomy.InCore.not_integral_neg}
\leanok
\uses{thm:continuum_we_core}
If preferences are a.e.\ transitive and strongly monotone and $\{a : y(a) \succ_a x(a)\}$ is null-measurable for every
a.e.-measurable $y$, no integrable selection $g$ of $G$ has $\int g \ll 0$.
\end{lemma}
\begin{proof}
\leanok
Let $S$ be (a measurable version of) the set where $g(a) \ne 0$ lies in the first part of $G(a)$; $\int_S g = \int g \ll 0$,
so $\mu(S) > 0$. Give each member of $S$ the bundle $e(a) + g(a) + w$ with $w = -\int g / \mu(S) \gg 0$. By strong
monotonicity and transitivity each member strictly prefers it to $x(a)$, and $\int_S$ of it equals $\int_S e$: $S$ blocks.
\end{proof}

\begin{lemma}[Supporting prices]
\label{lem:aumann_price}
\lean{Econ.Continuum.ContinuumEconomy.InCore.exists_price}
\leanok
\uses{lem:aumann_no_neg, thm:convex_selection_integral}
For finite nonatomic $\mu$ there is $p \ge 0$, $p \ne 0$, with $p \cdot \int g \ge 0$ for every integrable selection $g$ of
$G$.
\end{lemma}
\begin{proof}
\leanok
$\int G$ is convex (Theorem~\ref{thm:convex_selection_integral}) and misses the open negative orthant; Hahn--Banach separates
them by a functional $f$ and level $u$. Since $0 \in \int G$ and the orthant is a cone, $u = 0$, $f \le 0$ on the orthant,
so $f = p \cdot$ with $p \ge 0$, $\sum p > 0$.
\end{proof}

\begin{lemma}[Price support almost everywhere]
\label{lem:aumann_support}
\lean{Econ.Continuum.ContinuumEconomy.ae_price_support, Econ.Continuum.ContinuumEconomy.setIntegral_neg_of_ae_neg}
\leanok
\uses{lem:aumann_price}
If moreover the strict upper set of $x(a)$ is relatively open in $\mathbb{R}^\iota_+$ for a.e.\ $a$, then for a.e.\ $a$,
every $z \ge 0$ with $z \succ_a x(a)$ has $p \cdot z \ge p \cdot e(a)$.
\end{lemma}
\begin{proof}
\leanok
For a rational bundle $c \ge 0$, if the traders preferring $c$ to $x(a)$ with $p \cdot c < p \cdot e(a)$ had positive
measure, the selection $1_T (c - e)$ would have negative value. So each such set is null; there are countably many. If
$z \succ_a x(a)$ with $p \cdot z < p \cdot e(a)$, by continuity a rational $c$ slightly above $z$ has the same properties.
\end{proof}

\begin{theorem}[Aumann 1964: core $\subseteq$ Walrasian]
\label{thm:aumann_core}
\lean{Econ.Continuum.ContinuumEconomy.InCore.isWalrasianEquilibrium, Econ.Continuum.ContinuumEconomy.not_P_of_cost_le,
  Econ.Continuum.ContinuumEconomy.eq_zero_of_cost_nonpos, Econ.Continuum.ContinuumEconomy.cost_const_eq}
\leanok
\uses{lem:aumann_support}
Let $\mu$ be finite and nonatomic, with at least one good and $\int e \gg 0$, and let preferences be a.e.\ rational and
strongly monotone. If $x$ is in the core and satisfies the continuity and measurability conditions above, then $x$ is a
Walrasian allocation at some $p \gg 0$.
\end{theorem}
\begin{proof}
\leanok
Monotonicity and price support give $p \cdot x(a) \ge p \cdot e(a)$ a.e., and feasibility forces equality. Cheaper point: if
$p \cdot e(a) > 0$, a preferred affordable $y$ could be shrunk to $sy$, $s < 1$, still preferred by continuity but strictly
cheaper than $e(a)$, contradicting price support. Applied to $x(a) + \mathbf{1}_i$ for a trader with $p \cdot e(a) > 0$
(which exist as $\int e \gg 0$) this gives $p_i > 0$. Traders with $p \cdot e(a) = 0$ have $e(a) = x(a) = 0$ and can afford
only $0$.
\end{proof}

\begin{corollary}[Core equivalence]
\label{cor:aumann_equivalence}
\lean{Econ.Continuum.ContinuumEconomy.eventually_P_of_continuousPref, Econ.Continuum.ContinuumEconomy.nullMeasurableSet_P,
  Econ.Continuum.ContinuumEconomy.InCore.isWalrasianEquilibrium_of_continuousPref,
  Econ.Continuum.ContinuumEconomy.inCore_iff_of_continuousPref}
\leanok
\uses{thm:aumann_core}
With a.e.\ rational, strongly monotone, continuous preferences and jointly measurable strict preference
$\{(a,u,v) : u \succ_a v\}$, the core equals the set of Walrasian allocations. No convexity is assumed.
\end{corollary}
\begin{proof}
\leanok
Closed lower sets give relatively open strict upper sets; joint measurability composed with $a \mapsto (a, y(a), x(a))$
gives the measurability condition. The converse is Theorem~\ref{thm:continuum_we_core}.
\end{proof}

\begin{theorem}[Witnesses]
\label{thm:aumann_witness}
\lean{Econ.Continuum.ContinuumEconomy.linContinuum_inCore_iff, Econ.Continuum.ContinuumEconomy.uBump,
  Econ.Continuum.ContinuumEconomy.continuous_uBump_uncurry, Econ.Continuum.ContinuumEconomy.continuous_uBump,
  Econ.Continuum.ContinuumEconomy.stronglyMonotone_uBump, Econ.Continuum.ContinuumEconomy.bumpContinuum,
  Econ.Continuum.ContinuumEconomy.bumpContinuum_not_convexPref, Econ.Continuum.ContinuumEconomy.bumpContinuum_inCore_iff}
\leanok
\uses{cor:aumann_equivalence, thm:lyapunov_integral_witness}
Core equivalence holds for uniform traders on $[0,1]$ with endowment $(1,1)$ and utility $z_0 + z_1$, and also with the
non-convex, trader-dependent utility $\max(z_0, z_1) + (1 + a^2)(z_0 + z_1)$.
\end{theorem}
\begin{proof}
\leanok
All hypotheses are checked directly; for the second, $(1,0)$ and $(0,1)$ are indifferent but their midpoint is worse.
\end{proof}

\section{Existence in continuum economies}
\label{sec:continuum_existence}

\begin{lemma}[Splitting a population by prescribed masses]
\label{lem:split_masses}
\lean{Econ.Lyapunov.exists_partition_measureReal_eq_fin, Econ.Lyapunov.exists_partition_measureReal_eq}
\leanok
\uses{thm:sierpinski}
Let $\mu$ be finite and nonatomic, $A$ measurable, and $c_k \ge 0$ with $\sum_k c_k = \mu(A)$ over a nonempty finite index
set. Then $A$ is covered by pairwise disjoint measurable $B_k \subseteq A$ with $\mu(B_k) = c_k$.
\end{lemma}
\begin{proof}
\leanok
Induction on the number of indices, peeling off a subset of measure $c_0$ by Sierpi\'nski's theorem
(Theorem~\ref{thm:sierpinski}) and applying the inductive hypothesis to $A \setminus B_0$.
\end{proof}

\begin{lemma}[Carath\'eodory with a fixed index set]
\label{lem:caratheodory_fixed}
\lean{Econ.sum_extend_of_injective, Econ.exists_fin_convexComb}
\leanok
Every point of $\operatorname{conv} s$, for nonempty $s$ in a finite-dimensional real space $V$, is a convex combination of
$\dim V + 1$ points of $s$ (weights may vanish).
\end{lemma}
\begin{proof}
\leanok
Carath\'eodory gives an affinely independent family, whose cardinality is at most $\dim V + 1$; pad the weights with zeros
along an injection into $\mathrm{Fin}(\dim V + 1)$.
\end{proof}

\begin{theorem}[Convexification preserves a closed graph]
\label{thm:closed_graph_conv}
\lean{Econ.convexHull_nonempty_iff, Econ.isClosed_graph_convexHull, Econ.isClosed_graph_convexHull_on}
\leanok
\uses{lem:caratheodory_fixed}
If $\Phi$ has a closed graph and all its values lie in one compact $K$, then $x \mapsto \operatorname{conv}(\Phi x)$ has a
closed graph (also relative to a domain $s$, the form Kakutani's theorem uses); it has nonempty values exactly where
$\Phi$ does, and convex values always.
\end{theorem}
\begin{proof}
\leanok
Represent $v_n \in \operatorname{conv}(\Phi x_n)$ by Lemma~\ref{lem:caratheodory_fixed}; the weights and points live in the
compact set $\Delta \times K^{N}$, so a subsequence converges, the limit points lie in $\Phi x$ by the closed graph, and the
limit of the combinations is the combination of the limits.
\end{proof}

\begin{theorem}[Convexification bridge]
\label{thm:convexification_bridge}
\lean{Econ.Continuum.ContinuumEconomy.comp_type_eq_sum, Econ.Continuum.ContinuumEconomy.integrable_comp_type,
  Econ.Continuum.ContinuumEconomy.integral_comp_type, Econ.Continuum.ContinuumEconomy.integral_endow_eq_sum,
  Econ.Continuum.ContinuumEconomy.exists_allocation_isWalrasianEquilibrium_of_mem_convexHull}
\leanok
\uses{lem:split_masses}
Let a continuum economy have finitely many types, trader $a$ having the preference and endowment of type $\tau(a)$, with
$\mu$ finite and nonatomic. If at prices $p$ each type $t$ has $\bar x_t \in \operatorname{conv} (\text{demand}_t(p))$ and
$\sum_t \mu(\tau^{-1} t)\, \bar x_t = \sum_t \mu(\tau^{-1} t)\, \omega_t$, then $p$ supports a Walrasian equilibrium of the
continuum economy.
\end{theorem}
\begin{proof}
\leanok
Write $\bar x_t = \sum_k w_{tk} z_{tk}$ with $z_{tk}$ demanded (Lemma~\ref{lem:caratheodory_fixed}), split the traders of
type $t$ into groups of masses $w_{tk}\mu(\tau^{-1} t)$ (Lemma~\ref{lem:split_masses}), and give group $k$ the bundle
$z_{tk}$. Each trader demands her bundle, and the aggregate reproduces the convex combination.
\end{proof}

\begin{theorem}[Existence with finitely many types, strictly convex preferences]
\label{thm:continuum_existence_types}
\lean{Econ.Continuum.ContinuumEconomy.demand_scaled,
  Econ.Continuum.ContinuumEconomy.strictlyQuasiconcave_scaled,
  Econ.Continuum.ContinuumEconomy.stronglyMonotone_scaled,
  Econ.Continuum.ContinuumEconomy.exists_isWalrasianEquilibrium_of_types}
\leanok
\uses{thm:convexification_bridge}
A continuum economy with finitely many types of positive mass, each with a continuous, strictly quasiconcave utility
representing strongly monotone preferences and a nonzero endowment, and with $\int e \gg 0$, has a Walrasian equilibrium at
strictly positive prices in which all traders of a type consume the same bundle. This is a reduction to the finite
theorem; nonatomicity is not needed.
\end{theorem}
\begin{proof}
\leanok
The consumer with utility $x \mapsto u_t(m_t^{-1}x)$ and endowment $m_t \omega_t$ demands $m_t$ times what type $t$ demands,
so the mass-weighted clearing condition is ordinary clearing in the finite auxiliary economy over the types, where existence
is MWG 17.B.2 + 17.C.1; give every trader of type $t$ the unscaled bundle.
\end{proof}

\begin{theorem}[Witness]
\label{thm:continuum_existence_witness}
\lean{Econ.Continuum.ContinuumEconomy.sqrtType, Econ.Continuum.ContinuumEconomy.measurable_sqrtType,
  Econ.Continuum.ContinuumEconomy.sqrtOmega, Econ.Continuum.ContinuumEconomy.sqrtContinuum,
  Econ.Continuum.ContinuumEconomy.sqrtType_fiber_zero, Econ.Continuum.ContinuumEconomy.sqrtType_fiber_one,
  Econ.Continuum.ContinuumEconomy.sqrtContinuum_mass_zero, Econ.Continuum.ContinuumEconomy.sqrtContinuum_mass_one,
  Econ.Continuum.ContinuumEconomy.sqrtContinuum_mass_pos, Econ.Continuum.ContinuumEconomy.sqrtContinuum_endow_integral,
  Econ.Continuum.ContinuumEconomy.sqrtContinuum_exists_equilibrium}
\leanok
\uses{thm:continuum_existence_types}
Traders uniform on $[0,1]$, utility $\sqrt{x_0} + \sqrt{x_1}$, endowment $(1,2)$ on $[0,1/3]$ and $(2,1)$ on $(1/3,1]$: an
equilibrium exists at strictly positive prices.
\end{theorem}
\begin{proof}
\leanok
The types have masses $1/3$ and $2/3$, and the aggregate endowment is $(5/3, 4/3) \gg 0$.
\end{proof}

\begin{theorem}[The continuum is essential]
\label{thm:continuum_nonconvex_witness}
\lean{Econ.Continuum.ContinuumEconomy.uMax, Econ.Continuum.ContinuumEconomy.maxContinuum,
  Econ.Continuum.ContinuumEconomy.one2, Econ.Continuum.ContinuumEconomy.uMax_le_two,
  Econ.Continuum.ContinuumEconomy.mem_demand_uMax, Econ.Continuum.ContinuumEconomy.maxContinuum_const_not_demanded,
  Econ.Continuum.ContinuumEconomy.maxContinuum_no_symmetric_equilibrium,
  Econ.Continuum.ContinuumEconomy.maxContinuum_exists_equilibrium}
\leanok
\uses{thm:convexification_bridge, thm:lyapunov_integral_witness}
Identical traders uniform on $[0,1]$ with endowment $(1,1)$ and the non-convex utility $\max(x_0, x_1)$: no equilibrium at
positive prices gives almost everyone the same bundle (so no type-constant allocation works), yet at prices $(1,1)$ there is an equilibrium, in which half the traders
consume $(2,0)$ and half $(0,2)$.
\end{theorem}
\begin{proof}
\leanok
An allocation that is a.e.\ constant and feasible is a.e.\ $(1,1)$, but spending all wealth on the cheaper good gives more than $1$ of it. At prices
$(1,1)$ both $(2,0)$ and $(0,2)$ are demanded and $(1,1)$ is their midpoint; apply
Theorem~\ref{thm:convexification_bridge}.
\end{proof}

\section{Existence without convex preferences (Aumann 1966)}
\label{sec:aumann_existence}

\begin{lemma}[Truncated demand without convexity]
\label{lem:truncdemand_nonconvex}
\lean{Econ.AumannExistence.cost_single, Econ.AumannExistence.truncDemand_closed_nonempty,
  Econ.AumannExistence.truncDemand}
\leanok
For a continuous utility $u$, box size $c \ge 0$ and floor $m > 0$, the maximizers of $u$ over
$\{x \in [0,c]^L \mid p \cdot x \le \max(w, m)\}$ are nonempty and form a closed graph in $(p, w, x)$. No quasiconcavity
is assumed. The truncated demand of a consumer with endowment $\omega$ at prices $q$ is this set at $w = q\cdot\omega$.
\end{lemma}
\begin{proof}
\leanok
Berge's maximum theorem: the truncated budget correspondence is continuous with nonempty compact values.
\end{proof}

\begin{lemma}[Demand explodes near a free good, uniformly in the truncation]
\label{lem:boundary_uniform}
\lean{Econ.AumannExistence.eventually_lt_sum_of_mem_truncDemand}
\leanok
\uses{lem:truncdemand_nonconvex}
Let $u$ be continuous and strongly monotone, $m > 0$, and $p_i = 0$. For every $B$ there is a neighbourhood of $p$ such
that for all $q$ in it and all $c \ge B + 1$, every truncated demand $d$ at $q$ has $\sum_j d_j > B$.
\end{lemma}
\begin{proof}
\leanok
Let $K = \{d \ge 0 \mid \sum_j d_j \le B\}$, compact. If $d_0 \in K$ is affordable at $p$, then $y = s(d_0 + e_i)$ for some
$s < 1$ is strictly better than $d_0$, strictly cheaper than the wealth $\max(p\cdot\omega, m) > 0$ (good $i$ is free), and
lies in every box with $c \ge B + 1$; by continuity $y$ stays affordable and strictly better than $d$ for $(q, d)$ near
$(p, d_0)$, so such $d$ is not a truncated demand. If $d_0$ is unaffordable at $p$, nearby $d$ is unaffordable at nearby
$q$. Compactness of $K$ makes the neighbourhood uniform.
\end{proof}

\begin{lemma}[Uniform lower bound on prices]
\label{lem:price_floor}
\lean{Econ.AumannExistence.exists_pos_le_of_notMem}
\leanok
If $U$ is a neighbourhood of every boundary point of the price simplex $\Delta$, there is $\pi > 0$ with $q_i \ge \pi$ for
all $i$ and all $q \in \Delta \setminus U$.
\end{lemma}
\begin{proof}
\leanok
$\Delta \setminus \operatorname{int} U$ is compact and consists of strictly positive prices; each coordinate attains a
positive minimum on it.
\end{proof}

\begin{theorem}[Equilibrium of the convexified economy]
\label{thm:convexified_equilibrium}
\lean{Econ.AumannExistence.convex_cost_le, Econ.AumannExistence.convex_le_cost, Econ.AumannExistence.convex_lt_cost,
  Econ.AumannExistence.sum_eq_cost_one, Econ.AumannExistence.exists_convexifiedEquilibrium}
\leanok
\uses{lem:boundary_uniform, lem:price_floor, thm:closed_graph_conv, thm:kakutani}
Let finitely many types $t$ have masses $m_t > 0$, continuous strongly monotone utilities $u_t$ (not necessarily
quasiconcave) and endowments $\omega_t \gg 0$. Then there are prices $p \gg 0$ and $\bar x_t \in \operatorname{conv}
(\text{demand}_t(p, p\cdot\omega_t))$ with $\sum_t m_t \bar x_t = \sum_t m_t \omega_t$.
\end{theorem}
\begin{proof}
\leanok
Let $S = \sum_t m_t\omega_t$, fix a type $t_0$ and $B = (\sum_i S_i)/m_{t_0}$. Lemmas~\ref{lem:boundary_uniform}
and~\ref{lem:price_floor} give $\pi > 0$ such that any $q \in \Delta$ with some $q_i < \pi$ has, for every $c \ge B+1$,
all truncated demands of type $t_0$ of total quantity $> B$. Take $c = B + 1 + W/\pi$ with $W = \sum_t\sum_i \omega_{ti}$.
Apply Kakutani (Theorem~\ref{thm:kakutani}) on $\Delta \times [0,c]^{T\times L}$ to
$(p, x) \mapsto \arg\max_{q \in \Delta} q\cdot z(x) \times \prod_t \operatorname{conv} D_t(p)$, where
$z(x) = \sum_t m_t x_t - S$ and $D_t$ is truncated demand; the graph is closed by Theorem~\ref{thm:closed_graph_conv}.
At a fixed point, budgets give $p \cdot z \le 0$, so $z \le 0$ by the choice of $p$; then $\sum_i x_{t_0 i} \le B$, and
since $\operatorname{conv} D_{t_0}(p)$ lies in the convex set $\{\sum_j v_j > B\}$ otherwise, $p_i \ge \pi$ for all $i$.
Every affordable bundle then has coordinates at most $W/\pi < c$, so truncated demand is true demand; strong monotonicity
makes budgets bind, so $p\cdot z = 0$ and, as $p \gg 0$, $z = 0$.
\end{proof}

\begin{theorem}[Aumann's existence theorem, finitely many types]
\label{thm:aumann_existence}
\lean{Econ.Continuum.ContinuumEconomy.exists_isWalrasianEquilibrium_of_types_nonconvex}
\leanok
\uses{thm:convexified_equilibrium, thm:convexification_bridge}
A finite nonatomic continuum economy with finitely many types of positive mass, each with a continuous strongly monotone
utility (no convexity) and endowment $\gg 0$, has a Walrasian equilibrium at strictly positive prices.
\end{theorem}
\begin{proof}
\leanok
Theorem~\ref{thm:convexified_equilibrium} with masses $\mu(\tau^{-1}t)$, then Theorem~\ref{thm:convexification_bridge}.
\end{proof}

\begin{remark}[Scope relative to Aumann 1966]
This covers finitely many types (not a measurable preference map), preferences given by continuous utilities (no loss
for complete, transitive, continuous preferences, but not Schmeidler's incomplete ones), and endowments $\gg 0$ for every
type, whereas Aumann and Hildenbrand need only the aggregate endowment $\int e \gg 0$. Strong monotonicity is Aumann's
desirability assumption.
\end{remark}

\begin{theorem}[A non-convex economy whose equilibria all split identical traders]
\label{thm:aumann_existence_witness}
\lean{Econ.Continuum.ContinuumEconomy.uSq, Econ.Continuum.ContinuumEconomy.sqContinuum,
  Econ.Continuum.ContinuumEconomy.continuous_uSq, Econ.Continuum.ContinuumEconomy.stronglyMonotone_uSq,
  Econ.Continuum.ContinuumEconomy.not_quasiconcaveOn_uSq,
  Econ.Continuum.ContinuumEconomy.sqContinuum_one2_not_demanded,
  Econ.Continuum.ContinuumEconomy.sqContinuum_no_symmetric_equilibrium,
  Econ.Continuum.ContinuumEconomy.sqContinuum_exists_equilibrium,
  Econ.Continuum.ContinuumEconomy.sqContinuum_exists_asymmetric_equilibrium}
\leanok
\uses{thm:aumann_existence, thm:lyapunov_integral_witness}
Traders uniform on $[0,1]$ with endowment $(1,1)$ and utility $x_0^2 + x_1^2$, continuous and strongly monotone on
$\mathbb{R}^2_+$ but not quasiconcave. An equilibrium at positive prices exists, and in every equilibrium at positive
prices the allocation is not a.e.\ constant.
\end{theorem}
\begin{proof}
\leanok
Existence is Theorem~\ref{thm:aumann_existence}. An allocation that is a.e.\ constant and feasible forces the constant
to be $(1,1)$, but spending
all wealth $p_0 + p_1$ on the cheaper good gives at least $2$ units, utility at least $4 > 2$.
\end{proof}
